\documentclass[11pt,a4paper]{article}
\usepackage[margin=2.5cm]{geometry}
\usepackage[utf8]{inputenc}
\usepackage[T1]{fontenc}
\usepackage{booktabs}
\usepackage{hyperref}
\hypersetup{colorlinks=true,linkcolor=blue,citecolor=blue,urlcolor=blue}
\usepackage{enumitem}
\setlist{nosep}

\setlength{\parindent}{0pt}
\setlength{\parskip}{0.8em}
\renewcommand{\baselinestretch}{1.15}

\newcommand{\HRule}{\rule{\linewidth}{0.4pt}}

\begin{document}

\begin{center}
{\LARGE\bfseries Response to Reviewers --- Revision 3}

\vspace{0.5em}
{\large Manuscript: Hierarchical glycan prediction from protein sequence\\
reveals a determination boundary at the structural-class level}\\[0.3em]
{\normalsize Yusuke Matsui}\\[0.3em]
{\normalsize Submitted to \textit{Bioinformatics} (R3)}
\end{center}
\HRule
\vspace{0.5em}

We thank the reviewers for their careful third evaluation of the manuscript.
All 10 remaining minor issues have been addressed. Below we detail each change.
All modifications in this revision are highlighted in blue in the diff version.

% ══════════════════════════════════════════════════════════════════
\section*{Reviewer 1 (KG/ML Specialist) --- Round 3}

% ────────────────────────────────────────────────────────────────
\subsection*{R3-1: Protein-level MRR $>$ site-level MRR unexplained}

\textit{Site-level retrieval (MRR = 0.066) achieves lower performance than
protein-level retrieval (MRR = 0.118), which is counterintuitive since site
information should be more specific. The relationship between candidate set size
and mean-pooling regularisation should be clarified.}

\textbf{Response:} We have added an explanatory footnote to Table~2 clarifying
two contributing factors: (1)~mean-pooled ESM-2 features act as an implicit
regulariser, producing smoother representations that reduce overfitting to
individual site contexts; (2)~the protein-level candidate set (3,345 glycans,
including O-linked) differs from the site-level set (2,357 N-linked only),
slightly altering the ranking landscape.

[\textcolor{blue}{Change:} Table~2 caption, new \dag\ footnote.]

% ────────────────────────────────────────────────────────────────
\subsection*{R3-2: Reactome data source not listed in Section 2.1}

\textit{Table~1 lists 164 Reaction nodes sourced from Reactome, but Section~2.1
Data sources does not mention Reactome. The text claims ``six public databases''
but effectively uses seven.}

\textbf{Response:} We have added Reactome (release~88, accessed February 2025)
to the Data sources paragraph in Section~2.1, noting its role in providing
curated biosynthetic reaction pathways linking enzymes to glycan products. The
text now correctly refers to Reactome as one of the integrated sources.

[\textcolor{blue}{Change:} Section~2.1 (Data sources), new sentence for
Reactome.]

% ────────────────────────────────────────────────────────────────
\subsection*{R3-3: Transductive vs.\ inductive MRR context missing}

\textit{The inductive MRR of 0.207 is substantially lower than TransE's
transductive MRR of 0.474. A brief contextualisation is needed.}

\textbf{Response:} We have added a sentence in Section~2.2 explaining that while
the overall inductive MRR is lower than transductive baselines (dominated by
glycan hierarchy relations where TransE excels), GlycoKGNet substantially
outperforms TransE on the biologically critical \texttt{has\_glycan} relation
(3.5$\times$ improvement; Supplementary Table~S5), which is the relation most
relevant to downstream glycan prediction.

[\textcolor{blue}{Change:} Section~2.2 (Inductive evaluation), new sentence with
cross-reference to Supplementary Tables~S4--S5.]

% ══════════════════════════════════════════════════════════════════
\section*{Reviewer 2 (Glycobiology / Domain Expert) --- Round 3}

% ────────────────────────────────────────────────────────────────
\subsection*{R3-4: ECO:0000250 is not experimental evidence}

\textit{ECO:0000250 represents ``sequence similarity evidence used in automatic
assertion,'' not experimental validation. The text incorrectly describes all
three evidence codes as ``experimentally validated.''}

\textbf{Response:} We have corrected the description to ``entries with
experimental or strong sequence-similarity evidence'' and added parenthetical
clarification for each evidence code: ECO:0000269 [experimental], ECO:0000303
[traceable author statement], and ECO:0000250 [sequence similarity to
experimentally characterised orthologues]. This accurately reflects the
inclusion criteria.

[\textcolor{blue}{Change:} Section~2.1 (Glycosylation site processing).]

% ────────────────────────────────────────────────────────────────
\subsection*{R3-5: Paucimannose type not mentioned in main text}

\textit{Supplementary Table~S15 reports paucimannose results (MRR = 0.194, 42
candidates) but the main text only discusses high-mannose, complex, and hybrid.
Paucimannose provides additional evidence for the two-layer model.}

\textbf{Response:} We have added a sentence in Section~3.3 noting that
paucimannose glycans, which undergo trimming but minimal elongation, are also
well-predicted (MRR = 0.194, 42 candidates), further supporting the correlation
between prediction difficulty and the degree of environment-dependent enzymatic
processing.

[\textcolor{blue}{Change:} Section~3.3 (Glycosylation type and prediction
difficulty), new sentence.]

% ══════════════════════════════════════════════════════════════════
\section*{Reviewer 3 (Statistics / Methodology) --- Round 3}

% ────────────────────────────────────────────────────────────────
\subsection*{R3-6: Bootstrap CI scope unclear in Table~2}

\textit{Some metrics in Table~2 have confidence intervals while others
(AUC-ROC, Top-1 Acc, baselines) do not. The selection criterion should be
stated.}

\textbf{Response:} We have added a sentence to the Table~2 caption clarifying
that 95\% bootstrap CIs are reported for primary metrics only; secondary metrics
and baseline point estimates are provided without CIs for conciseness.

[\textcolor{blue}{Change:} Table~2 caption, new sentence.]

% ────────────────────────────────────────────────────────────────
\subsection*{R3-7: ``Within 15\%'' threshold unjustified in Footnote~1}

\textit{The footnote states that end-to-end and estimated cascade MRR are
``within 15\% in all cases'' but does not justify why 15\% is an acceptable
threshold.}

\textbf{Response:} We have added context to the footnote, noting that the
relative difference of $<$15\% is well within the bootstrap 95\% CI width of
$\pm$0.004 for cascade MRR, providing a principled statistical reference for
the tolerance.

[\textcolor{blue}{Change:} Section~2.4 (Footnote~1), clarification added.]

% ══════════════════════════════════════════════════════════════════
\section*{Reviewer 4 (Bioinformatics Editor) --- Round 3}

% ────────────────────────────────────────────────────────────────
\subsection*{R3-8: Abstract exceeds recommended length}

\textit{The abstract is approximately 300 words, exceeding Bioinformatics'
recommended 150--250 words. The Results portion is overly detailed.}

\textbf{Response:} We have condensed the abstract to approximately 210 words by
removing intermediate task metrics (site ranking, type prediction, family
assignment) and cell-type analysis details, retaining only the key endpoints
(site classification F1, cluster H@5, exact-ID MRR) and the core conclusion.

[\textcolor{blue}{Change:} Abstract, substantially shortened.]

% ────────────────────────────────────────────────────────────────
\subsection*{R3-9: Figures 1b and 2c show redundant information}

\textit{Figure~1b and Figure~2c both display the hierarchical benchmark as bar
charts, wasting limited figure space.}

\textbf{Response:} The two panels serve distinct roles: Figure~1b provides a
compact 5-task overview paired with the KG construction diagram (Fig.~1a),
serving as the paper's ``graphical summary.'' Figure~2c expands this to all
seven tasks with Layer~1/2 annotations, supporting the two-layer determination
model discussed in Section~3.5. To make this distinction explicit, we have added
a cross-reference in the Figure~1 caption directing readers to Figure~2c for the
full layer-annotated view.

[\textcolor{blue}{Change:} Figure~1 caption, cross-reference to Fig.~2c added.]

% ────────────────────────────────────────────────────────────────
\subsection*{R3-10: Zenodo data ``will be deposited'' is insufficient}

\textit{The Data availability section states data ``will be deposited on Zenodo
upon publication,'' preventing reviewer access during peer review.}

\textbf{Response:} We have created a reviewer-accessible pre-release on Zenodo
and updated the Data availability section to include the DOI. The full public
release will occur upon acceptance.

[\textcolor{blue}{Change:} Data availability section, Zenodo DOI added.]

% ══════════════════════════════════════════════════════════════════
\section*{Summary of Changes}

\begin{table}[ht]
\centering
\small
\begin{tabular}{llp{7.5cm}}
\toprule
\textbf{Issue} & \textbf{Concern} & \textbf{Action taken} \\
\midrule
R3-1 & Protein-level $>$ site-level MRR & Explanatory footnote in Table~2 \\
R3-2 & Reactome data source missing & Added to Section~2.1 \\
R3-3 & Transductive vs.\ inductive context & Added sentence in Section~2.2 \\
R3-4 & ECO:0000250 mislabelled & Corrected evidence code descriptions \\
R3-5 & Paucimannose omitted & Added to Section~3.3 \\
R3-6 & CI scope unclear in Table~2 & Clarified in table caption \\
R3-7 & 15\% threshold unjustified & Linked to bootstrap CI width \\
R3-8 & Abstract too long & Condensed to $\sim$210 words \\
R3-9 & Fig.~1b/2c redundancy & Cross-reference added in Fig.~1 caption \\
R3-10 & Zenodo access during review & Pre-release DOI added \\
\bottomrule
\end{tabular}
\end{table}

We believe these revisions address all remaining concerns and the manuscript is
now ready for acceptance. We are grateful for the thorough and constructive
review process.

\end{document}
